\section{循环检测与计划重审}

\subsection{Doom Loop 问题}

Agent 一旦确定了某个方案，往往会表现出严重的\textbf{近视行为}（Myopic Behavior）：即使方案明显不可行，仍然反复对同一个文件做微小变种修改，形成所谓的"Doom Loop"。在 Trace 分析中，LangChain 观察到有些 Agent 在同一个错误方案上尝试了超过 10 次微调。

\begin{warningbox}{Doom Loop 的本质}
Doom Loop 反映了模型的一个深层问题：当模型对自己的方案产生了"信念锁定"，它会倾向于在\textbf{执行层面}做微调（调参数、改小细节），而不是在\textbf{策略层面}重新审视方案是否正确。这类似于人类的"沉没成本谬误"。
\end{warningbox}

\subsection{LoopDetectionMiddleware}

LangChain 设计了 LoopDetectionMiddleware 来应对这一问题：

\begin{itemize}[leftmargin=1.5em]
    \item \textbf{监控机制}：通过 Tool Call Hook 追踪每个文件的编辑次数
    \item \textbf{触发条件}：当同一文件的编辑次数超过阈值 $N$ 时触发
    \item \textbf{干预方式}：向 Agent 注入上下文信息，如"你已经对这个文件修改了多次，请考虑重新审视你的整体方案"
\end{itemize}

需要注意的是，这种干预是\textbf{建议性}的，而非强制性的——如果模型仍然认为当前方向正确，它可能会继续沿着原路走。

\begin{knowledgebox}{Harness 护栏的时效性}
LangChain 团队特别强调：LoopDetectionMiddleware 是一种针对\textbf{当前模型缺陷}的设计启发式（Design Heuristic）。随着模型能力的提升，这类护栏很可能会变得不再必要。但在当下，它们是帮助 Agent 正确、自主执行的有效工具。Harness 工程师需要\textbf{围绕当前模型的短板做设计，同时为更强模型做好演进规划}。
\end{knowledgebox}

\subsection{本章小结}

Doom Loop 是当前 Agent 最顽固的失败模式之一。LoopDetectionMiddleware 通过追踪文件编辑频率并注入重审提示来缓解这一问题。这类"护栏"设计具有时效性——它们是为当前模型的短板量身定制的临时方案，随着模型进化将逐渐退役。


